\documentclass[10pt]{article}
\setlength\parindent{10pt}

\title{\vspace{-2cm} \flushleft \normalsize Scientometrics\\[1cm]
\large\bfseries\centering Dual ordinal spectral ranking of journals and institutions}
\date{\today}
\author{\textsc{Docampo, D. and Cram, L.}}

\usepackage{natbib, amsmath, amssymb, mathrsfs}
\usepackage[margin=1in]{geometry}
\usepackage{tikz}
\usetikzlibrary{positioning, arrows.meta}
\usepackage{hyperref}
\pagestyle{plain}

\usepackage{filecontents}
\begin{filecontents}{birank_assembled.bib}
@phdthesis{woutersThesisCitationCulture1999,
  author = {Wouters, Paul},
  title  = {The Citation Culture},
  school = {University of Amsterdam},
  year   = {1999}
}
@article{pinskiNarin1976,
  author  = {Pinski, Gabriel and Narin, Francis},
  title   = {Citation influence for journal aggregates of scientific publications},
  journal = {Information Processing \& Management},
  volume  = {12},
  pages   = {297--312},
  year    = {1976}
}
@article{geller1978,
  author  = {Geller, Norman L.},
  title   = {On the citation influence methodology of {Pinski} and {Narin}},
  journal = {Information Processing \& Management},
  volume  = {14},
  pages   = {93--95},
  year    = {1978}
}
@article{palacioshuertaVolij2004,
  author  = {Palacios-Huerta, Ignacio and Volij, Oscar},
  title   = {The measurement of intellectual influence},
  journal = {Econometrica},
  volume  = {72},
  pages   = {963--977},
  year    = {2004}
}
@article{breigerDualityPersonsGroups1974a,
  author  = {Breiger, Ronald L.},
  title   = {The duality of persons and groups},
  journal = {Social Forces},
  volume  = {53},
  pages   = {181--190},
  year    = {1974}
}
@article{katz1953,
  author  = {Katz, Leo},
  title   = {A new status index derived from sociometric analysis},
  journal = {Psychometrika},
  volume  = {18},
  pages   = {39--43},
  year    = {1953}
}
@article{franceschetPagerank2010,
  author  = {Franceschet, Massimo},
  title   = {{PageRank}: Standing on the shoulders of giants},
  journal = {Communications of the {ACM}},
  volume  = {54},
  pages   = {92--101},
  year    = {2011}
}
@article{vignaSpectralRanking2016,
  author  = {Vigna, Sebastiano},
  title   = {Spectral ranking},
  journal = {Network Science},
  volume  = {4},
  pages   = {433--445},
  year    = {2016}
}
@article{priemOpenAlexFullyopenIndex2022,
  author  = {Priem, Jason and Piwowar, Heather and Orr, Richard},
  title   = {{OpenAlex}: A fully-open index of scholarly works, authors, venues, institutions, and concepts},
  journal = {arXiv preprint arXiv:2205.01833},
  year    = {2022}
}
\end{filecontents}

\begin{document}
\maketitle

\begin{abstract}
We develop a dual spectral ranking of journals and institutions that
jointly determines prestige scores via a Katz resolvent applied to a
block attention matrix constructed from intra-field citation flows.
A continuous parameter $\chi$ apportions each reference between source
and institution channels; a binary block-mask selects among four
bibliometrically distinct flow structures, nesting conventional
single-mode rankings as special cases.
Applied to the economics and commerce literature using OpenAlex data,
the method produces stable, interpretable joint rankings whose
sensitivity to model parameters is examined systematically.

\medskip\noindent
\textbf{Keywords}: Scientometrics, spectral ranking, dual
journal--institution prestige, Katz resolvent, Pinski--Narin
\end{abstract}

% ---------------------------------------------------------------
\section{Introduction}
% ---------------------------------------------------------------

\paragraph{Journal and institution ranking.}
[Placeholder.]

\paragraph{Spectral methods.}
\citet{pinskiNarin1976} introduced influence-per-reference weights for
journal ranking; \citet{vignaSpectralRanking2016} surveys the broader
lineage; \citet{franceschetPagerank2010} traces the
Katz--Hubbell--Pinski-Narin--PageRank chain.
Institutional rankings (QS, THE, ARWU, Leiden) mix bibliometrics with
reputational surveys; spectral methods have not previously been applied
jointly to sources and institutions.

\paragraph{Theoretical framing.}
We treat citations as structural relations within a functionally
differentiated scientific system \citep{woutersThesisCitationCulture1999},
setting aside authorial intent.
The resulting dual rankings are self-referential hierarchies that
condition future citations.
Our construction extends \citet{breigerDualityPersonsGroups1974a}:
all four blocks of the joint attention matrix are primary flows, not
derived projections.

\paragraph{Paper structure.}
Section~\ref{sec:networks} develops the framework.
Section~\ref{sec:empirical} describes the corpus.
Section~\ref{sec:results} presents results.
Section~\ref{sec:discussion} discusses validity.

% ---------------------------------------------------------------
\section{Journal--institution citation networks}
\label{sec:networks}
% ---------------------------------------------------------------

\paragraph{Corpus construction.}

Following \citet{woutersThesisCitationCulture1999}, a \emph{reference}
written by a scientist becomes a \emph{citation}---a directed edge
$(w_i,w_j)$---only after indexer resolution to canonical work records.
The corpus is built in five steps:
(i)~select $N_s$ sources $\mathfrak{S}$ by field classification;
(ii)~collect all works $\mathfrak{W}$ ($|\mathfrak{W}|=N_w$) in
$\mathfrak{S}$ within the census window $[t^c_0,t^c_1]$;
(iii)~retain references from $\mathfrak{W}$ to works $\mathfrak{W}^t$
in $\mathfrak{S}$ within the target window $[t^t_0,t^t_1]$, forming
$\mathcal{R}=\{(i,j):w_i\in\mathfrak{W},\,w_j\in\mathfrak{W}^t\}$;
(iv)~induce institution set $\mathfrak{I}'$ from author affiliations
of $\mathfrak{W}$;
(v)~retain institutions with $\geq\tau$ census works, giving
$\mathfrak{I}$, $|\mathfrak{I}|=N_u$.
Flows crossing any corpus boundary are excluded without renormalisation
(\emph{open-corpus principle}); the closed-corpus alternative is
examined in the Supplement.

Each work $w_i$ is assigned an institution weight under
author-fractional counting,
\begin{equation}
  \omega_{ik} \;=\;
  \sum_{j:\,k\in\mathcal{U}_{ij}}
  \frac{1}{A_i}\,\frac{1}{U_{ij}},
  \label{eq:omega}
\end{equation}
where $A_i$ is the number of authors and $U_{ij}$ the number of
retained institutions of author $j$; after filtering,
$\sum_k\omega_{ik}\leq 1$ and weights are not renormalised.
Model parameters are
$\boldsymbol{\theta}=(\mathfrak{S},[t^c_0,t^c_1],[t^t_0,t^t_1],
\tau,\mathbf{m},\chi,\alpha)$.

\paragraph{The attention matrix $\mathbf{C}(\chi,\mathbf{m})$.}

A reference $(i,j)\in\mathcal{R}$ transfers one unit of attention from
$w_i$ to $w_j$.
On each side, a fraction $(1-\chi)$ is attributed to the source and
$\chi$ to the affiliated institutions (weighted by $\omega_{ik}$), so
each reference generates four typed edges with weights
$(1-\chi)^2$, $\chi(1-\chi)$, $\chi(1-\chi)$, $\chi^2$ for the
$SS$, $SI$, $IS$, $II$ blocks respectively.
Summing over $\mathcal{R}$ and applying a binary block-mask
$\mathbf{m}=(m_{SS},m_{SI},m_{IS},m_{II})\in\{0,1\}^4$,
with $m_{SI}=m_{IS}$, gives
\begin{equation}
  \mathbf{C}(\chi,\mathbf{m}) =
  \begin{pmatrix}
    m_{SS}(1-\chi)^2\,\mathbf{C}_{SS} &
    m_{SI}\,\chi(1-\chi)\,\mathbf{C}_{SI} \\[4pt]
    m_{IS}\,\chi(1-\chi)\,\mathbf{C}_{IS} &
    m_{II}\,\chi^2\,\mathbf{C}_{II}
  \end{pmatrix}.
  \label{eq:C}
\end{equation}
Four masks are bibliometrically meaningful: source-only $(1,0,0,0)$,
institution-only $(0,0,0,1)$, bipartite $(0,1,1,0)$, and full-joint
$(1,1,1,1)$.
The parameter $\chi$ controls the source--institution balance in the
bipartite and full-joint cases and is inert in the single-mode cases.
Block construction details are in the Supplement.

\paragraph{Spectral ranking.}

Let $\mathcal{P}=\mathfrak{S}\cup\mathfrak{I}$ be the \emph{unit} set
\citep{pinskiNarin1976}, $N=N_s+N_u$, $r_p=\sum_q C_{pq}$ the
outgoing reference flow, $\mathbf{D}_r=\mathrm{diag}(r_p)$, and
$\mathbf{A}=\mathrm{diag}(a_p)$ the work-count matrix ($a_p=N^s_p$
for sources, $a_p=M_k$ for institutions).
The row-stochastic reference matrix is
$\mathbf{H}=\mathbf{D}_r^{-1}\mathbf{C}$.
The Katz resolvent \citep{katz1953,franceschetPagerank2010},
\begin{equation}
  \boldsymbol{\pi}^*(\alpha)
  \;=\;(1-\alpha)\bigl(\mathbf{I}-\alpha\mathbf{H}^\top\bigr)^{-1}\mathbf{p},
  \quad p_p=\frac{a_p}{N_\text{works}},
  \label{eq:katz}
\end{equation}
converges for all $\alpha<1$ and approaches the stationary distribution
of the Markov chain with transition matrix $\mathbf{H}$ as $\alpha\to
1$; \citet{geller1978} shows $\pi^*_q$ equals the long-run expected
share of reference flow received by unit $q$.
\emph{Prestige per work} \citep{palacioshuertaVolij2004} is
\begin{equation}
  \mathbf{v}^*(\alpha)\;=\;\mathbf{A}^{-1}\boldsymbol{\pi}^*(\alpha),
  \label{eq:vstar}
\end{equation}
recovered once after convergence of successive substitution on
\eqref{eq:katz}. Convergence analysis is in the Supplement.
We report $\mathbf{v}^*$ over $\alpha\in\{0.50,0.75,0.85,0.95\}$.

\paragraph{Computation.}

$\mathbf{C}(\chi,\mathbf{m})$ is sparse ($\mathrm{nnz}\leq
4|\mathcal{R}|$) and is built from a weighted edge list $\mathcal{E}$
accumulated over $\mathcal{R}$, stored in CSR format, and
row-normalised to $\mathbf{H}$.
The pipeline
$\mathcal{R}\!\to\!\mathcal{E}\!\to\!\mathbf{C}\!\to\!\mathbf{H}\!\to\!\boldsymbol{\pi}^*\!\to\!\mathbf{v}^*$
and its parameter entry points are detailed in the Supplement.

% ---------------------------------------------------------------
\section{Example: Economics and Commerce}
\label{sec:empirical}
% ---------------------------------------------------------------

\paragraph{Data.}
\label{sec:data}
Guided by InCites EC journal categories, we extract OpenAlex data
\citep{priemOpenAlexFullyopenIndex2022} for approximately 600 journals
and 350,000 articles and reviews published 2000--2025.
Network properties and data limitations (reference cross-linking,
author disambiguation, institution identity) are detailed in the
Supplement.

% ---------------------------------------------------------------
\section{Results}
\label{sec:results}
% ---------------------------------------------------------------

\paragraph{Baseline model.} [Placeholder.]

\paragraph{Parameter sensitivity.}
[Placeholder: $\alpha\in\{0.50,0.75,0.85,0.95\}$;
$\chi\in\{0.00,0.25,0.50,0.75,1.00\}$;
$\mathbf{m}\in\{(1,0,0,0),(0,1,1,0),(1,1,1,1)\}$; $\tau$ sweep.]

\paragraph{Bootstraps.} [Placeholder.]

% ---------------------------------------------------------------
\section{Discussion}
\label{sec:discussion}
% ---------------------------------------------------------------

\paragraph{Validity.} [Placeholder: face, content, criterion, construct.]

% ---------------------------------------------------------------
\section{Prospects}
% ---------------------------------------------------------------

[Placeholder.]

\bibliographystyle{plainnat}
\bibliography{birank_assembled}

\end{document}
